% ===== Topic 1, part d (VI): problems 41-50 (nâng cao, mở rộng ý tưởng các bài 16-30) =====
% helper shapes for problem 41 (own prefix \rvab..., independent of part b)
\providecommand{\rvabShp}[3]{%
  \if#1T\draw[tmBlue!80!black,line width=0.8pt,fill=tmBlue!18] (#2,#3)--(#2+0.34,#3)--(#2+0.17,#3+0.3)--cycle;\fi
  \if#1C\draw[tmOrange!85!black,line width=0.8pt,fill=tmOrange!22] (#2+0.15,#3+0.15) circle (0.15);\fi
  \if#1S\draw[tmGreen!80!black,line width=0.8pt,fill=tmGreen!18] (#2,#3) rectangle (#2+0.3,#3+0.3);\fi
  \if#1D\draw[tmPurple!85!black,line width=0.8pt,fill=tmPurple!18] (#2+0.13,#3)--(#2+0.26,#3+0.16)--(#2+0.13,#3+0.32)--(#2,#3+0.16)--cycle;\fi
  \if#1Q\node[font=\Large\bfseries,text=Ink,inner sep=0pt,anchor=south] at (#2+0.15,#3) {?};\fi}
\providecommand{\rvabTri}{\tikz[baseline=0pt]{\rvabShp{T}{0}{0}}}
\providecommand{\rvabCir}{\tikz[baseline=0pt]{\rvabShp{C}{0}{0}}}
\providecommand{\rvabSq}{\tikz[baseline=0pt]{\rvabShp{S}{0}{0}}}
\providecommand{\rvabDia}{\tikz[baseline=0pt]{\rvabShp{D}{0}{0}}}
\providecommand{\rvabRow}[3]{%
  \foreach \rvabS [count=\rvabI] in {#1}{\xdef\rvabN{\rvabI}}%
  \foreach \rvabS [count=\rvabI] in {#1}{\rvabShp{\rvabS}{#2-0.19*\rvabN+0.38*\rvabI-0.36}{#3+0.03}}}
\providecommand{\rvabScale}[4]{\begin{scope}[shift={(#1,#2)}]
  \draw[Ink,line width=1pt] (-1.9,0.55)--(-0.5,0.55) (-1.2,0.55)--(-1.2,0.25)--(1.2,0.25)--(1.2,0.55) (0.5,0.55)--(1.9,0.55);
  \draw[Ink,line width=0.9pt,fill=Ink!12] (-0.2,-0.12)--(0,0.25)--(0.2,-0.12)--cycle;
  \rvabRow{#3}{-1.2}{0.55}\rvabRow{#4}{1.2}{0.55}
\end{scope}}

% ------------------------------------------------------------------ 41
\begin{probbox}{41}{\stars{1}}
Ba chiếc cân dưới đây đều thăng bằng; các hình giống nhau có khối lượng bằng nhau.
\begin{figC}
\rvabScale{0}{0}{T,T}{C,C,C}
\rvabScale{4.4}{0}{C,S}{T}
\rvabScale{8.8}{0}{D}{T,S}
\node[font=\footnotesize\bfseries,text=Slate] at (0,-0.45) {Cân 1};
\node[font=\footnotesize\bfseries,text=Slate] at (4.4,-0.45) {Cân 2};
\node[font=\footnotesize\bfseries,text=Slate] at (8.8,-0.45) {Cân 3};
\end{figC}
\begin{enumerate}[label=\alph*),leftmargin=1.6em,itemsep=1pt]
\item Cần bao nhiêu hình \rvabSq\ để cân bằng với một hình \rvabCir?
\item Cần bao nhiêu hình \rvabSq\ để cân bằng với một hình \rvabDia?
\item Đặt một \rvabDia\ và một \rvabCir\ lên đĩa trái. Nếu đĩa phải chỉ được dùng các hình \rvabTri\ thì cần bao nhiêu hình \rvabTri?
\item Biết tổng khối lượng của bốn vật \rvabTri, \rvabCir, \rvabSq, \rvabDia\ (mỗi loại một vật) là $50$~g. Tìm khối lượng của mỗi vật.
\end{enumerate}
\end{probbox}
\probsol{
\textbf{Phân tích đề.} Đây là bài mở rộng của Bài 17. Gọi $T$, $C$, $S$, $D$ (gam) lần lượt là khối lượng của một hình \rvabTri, một hình \rvabCir, một hình \rvabSq\ và một hình \rvabDia. Mỗi cân thăng bằng cho ta một đẳng thức (khối lượng đĩa trái $=$ khối lượng đĩa phải):
\[ \text{(1) } 2T=3C,\qquad \text{(2) } C+S=T,\qquad \text{(3) } D=T+S. \]
Ý tưởng: biểu diễn mọi khối lượng theo cùng một ẩn là $S$.
\begin{recallbox}Phương pháp thế: nếu biết $T=C+S$ thì ở bất kỳ đẳng thức nào ta cũng có thể thay $T$ bởi $C+S$. Cộng hoặc trừ cùng một biểu thức ở hai vế của một đẳng thức, ta được một đẳng thức đúng (quy tắc chuyển vế).\end{recallbox}

\step{1}{Câu a -- Thế (2) vào (1)}
Thay $T=C+S$ vào (1): $2(C+S)=3C$, tức $2C+2S=3C$.\par
Trừ $2C$ ở cả hai vế: $2S=3C-2C=C$. Vậy $C=2S$.\par
Một hình \rvabCir\ nặng bằng hai hình \rvabSq: cần \textbf{$2$ hình \rvabSq}.
\begin{tipbox}Vì sao lại thế (2) vào (1)? Vì (1) chỉ chứa $T$ và $C$; sau khi thay $T=C+S$, đẳng thức chỉ còn $C$ và $S$, nên so sánh được ngay $C$ với $S$.\end{tipbox}

\step{2}{Biểu diễn $T$ theo $S$}
Từ (2): $T=C+S=2S+S=3S$.\par
(Kiểm tra với (1): $2T=6S$ và $3C=3\cdot 2S=6S$ \checkmark.)

\step{3}{Câu b}
Từ (3): $D=T+S=3S+S=4S$. Cần \textbf{$4$ hình \rvabSq}.

\step{4}{Câu c}
Đĩa trái: $D+C=4S+2S=6S$. Mỗi \rvabTri\ nặng $T=3S$, nên số hình \rvabTri\ cần dùng là $6S:3S=$ \textbf{$2$ hình \rvabTri}. (Thật vậy, $2T=6S=D+C$.)

\step{5}{Câu d -- Lập phương trình một ẩn}
$T+C+S+D=50$, thay theo $S$: $3S+2S+S+4S=50$, tức $10S=50$, nên $S=5$. Suy ra:
\begin{center}
\begin{tabular}{|c|c|c|c|c|}
\hline
\rowcolor{tmBlue!10}Vật & \rvabSq & \rvabCir & \rvabTri & \rvabDia\\ \hline
Theo $S$ & $S$ & $C=2S$ & $T=3S$ & $D=4S$\\ \hline
Khối lượng & $5$~g & $10$~g & $15$~g & $20$~g\\ \hline
\end{tabular}
\end{center}
\verify{Cân 1: $2\cdot 15=30=3\cdot 10$ \checkmark. Cân 2: $10+5=15$ \checkmark. Cân 3: $20=15+5$ \checkmark. Tổng: $15+10+5+20=50$ \checkmark.}
\begin{parentbox}Hãy cho con đọc mỗi chiếc cân thành một đẳng thức bằng chữ ($2T=3C$,\dots) trước khi tính. Mỗi lần ``bỏ cùng một vật ở hai đĩa'' chính là trừ cùng một số ở hai vế -- đây là nền tảng của việc giải phương trình.\end{parentbox}
\ans{a) $2$ hình \rvabSq;\ b) $4$ hình \rvabSq;\ c) $2$ hình \rvabTri;\ d) \rvabSq\ $5$~g, \rvabCir\ $10$~g, \rvabTri\ $15$~g, \rvabDia\ $20$~g.}
}

% ------------------------------------------------------------------ 42
\begin{probbox}{42}{\stars{1}}
Trong bài này ta chỉ xét các tam giác mà số đo mỗi góc (tính bằng độ) là một \emph{số nguyên}. Hai tam giác cân được coi là \emph{cùng kiểu} nếu chúng có cùng bộ ba số đo góc.
\begin{enumerate}[label=\alph*),leftmargin=1.6em,itemsep=1pt]
\item Chứng minh rằng không có tam giác nào mà cả ba góc đều có số đo lẻ.
\item Có bao nhiêu kiểu tam giác cân mà cả ba góc đều có số đo chẵn?
\item Có bao nhiêu kiểu tam giác cân có đúng hai góc có số đo lẻ? Có tam giác cân nào có đúng một góc có số đo lẻ không?
\item Có tất cả bao nhiêu kiểu tam giác cân?
\end{enumerate}
\end{probbox}
\probsol{
\textbf{Phân tích đề.} Mở rộng của Bài 20: thay vì chỉ xét tam giác có góc $70^\circ$, ta xét \emph{mọi} tam giác cân có góc nguyên và phân loại chúng theo tính chẵn -- lẻ của các góc.
\begin{recallbox}Tổng ba góc của một tam giác bằng $180^\circ$. Tam giác cân có hai góc ở đáy bằng nhau. Về chẵn -- lẻ: lẻ $+$ lẻ $=$ chẵn; chẵn $+$ lẻ $=$ lẻ; tổng của ba số lẻ là số lẻ; $2\times$(số bất kỳ) là số chẵn.\end{recallbox}

\step{1}{Câu a}
Giả sử cả ba góc đều lẻ. Tổng hai số lẻ đầu là số chẵn; cộng thêm một số lẻ nữa được số \emph{lẻ}. Nhưng tổng ba góc là $180$ -- số chẵn. Mâu thuẫn. Vậy không có tam giác nào có cả ba góc lẻ.

\step{2}{Mô tả mọi tam giác cân bằng một số $x$}
Gọi số đo mỗi góc ở đáy là $x$ (độ). Khi đó góc ở đỉnh là $180-2x$. Điều kiện để là tam giác: góc ở đỉnh phải dương, tức $180-2x\ge 1$, hay $x\le 89$. Vậy $x$ có thể là $1,2,3,\dots,89$, và mỗi giá trị của $x$ cho đúng một kiểu tam giác cân.
\begin{proofbox}Hai giá trị $x$ khác nhau cho hai bộ góc khác nhau (vì góc ở đáy khác nhau), và mọi tam giác cân đều có dạng $(x;x;180-2x)$. Tam giác đều $(60;60;60)$ ứng với $x=60$ -- chỉ được đếm một lần. Lưu ý: ta luôn có thể coi góc ở đáy là hai góc bằng nhau; với tam giác đều thì góc nào cũng được.\end{proofbox}

\step{3}{Góc ở đỉnh luôn chẵn}
$180-2x=2\cdot(90-x)$ là một số chẵn, với mọi $x$. Do đó tính chẵn -- lẻ của tam giác chỉ phụ thuộc vào $x$:
\begin{itemize}[leftmargin=1.3em,itemsep=1pt]
\item $x$ chẵn: hai góc đáy chẵn, góc đỉnh chẵn $\Rightarrow$ cả ba góc chẵn.
\item $x$ lẻ: hai góc đáy lẻ, góc đỉnh chẵn $\Rightarrow$ đúng hai góc lẻ.
\end{itemize}

\step{4}{Câu b}
Đếm số chẵn từ $1$ đến $89$: $2,4,\dots,88$, có $88:2=44$ số. Vậy có \textbf{$44$ kiểu}.

\step{5}{Câu c}
Đếm số lẻ từ $1$ đến $89$: $1,3,\dots,89$, có $(89+1):2=45$ số. Vậy có \textbf{$45$ kiểu} có đúng hai góc lẻ.\par
Tam giác cân có đúng một góc lẻ thì \textbf{không có}: theo Bước 3, số góc lẻ chỉ có thể là $0$ hoặc $2$.

\step{6}{Câu d}
Tổng số kiểu: $44+45=89$ (đúng bằng số giá trị của $x$).
\begin{center}\small
\begin{tabular}{|c|c|c|}
\hline
\rowcolor{tmBlue!10}$x$ (góc đáy) & Bộ ba góc & Số góc lẻ\\ \hline
$1$ & $1^\circ;\,1^\circ;\,178^\circ$ & $2$\\ \hline
$2$ & $2^\circ;\,2^\circ;\,176^\circ$ & $0$\\ \hline
$55$ & $55^\circ;\,55^\circ;\,70^\circ$ & $2$\\ \hline
$70$ & $70^\circ;\,70^\circ;\,40^\circ$ & $0$\\ \hline
$89$ & $89^\circ;\,89^\circ;\,2^\circ$ & $2$\\ \hline
\end{tabular}
\end{center}
\begin{warnbox}Đừng quên các tam giác ``rất dẹt'' như $(1^\circ;1^\circ;178^\circ)$ -- chúng vẫn là tam giác hợp lệ. Cũng đừng cho $x=90$: khi đó góc đỉnh bằng $0^\circ$, không còn là tam giác.\end{warnbox}
\verify{Liệt kê bằng máy tính $x=1,\dots,89$: đúng $44$ bộ toàn chẵn và $45$ bộ có hai góc lẻ. Hai dòng $x=55$ và $x=70$ chính là hai tam giác của Bài 20.}
\ans{a) Không tồn tại;\ b) $44$;\ c) $45$ kiểu, và không có tam giác cân nào có đúng một góc lẻ;\ d) $89$.}
}

% ------------------------------------------------------------------ 43
\begin{probbox}{43}{\stars{2}}
Một chấm tròn xuất phát tại gốc tọa độ $O(0;0)$ và di chuyển theo luật như Bài 22: bước thứ nhất dài $1$ đơn vị theo phương ngang hoặc dọc; mỗi bước sau dài hơn bước trước $1$ đơn vị và vuông góc với bước ngay trước nó (chiều đi tùy ý).
\begin{enumerate}[label=\alph*),leftmargin=1.6em,itemsep=1pt]
\item Giải thích vì sao sau đúng $10$ bước, chấm không thể ở lại $O$.
\item Tìm số $n$ nhỏ nhất ($n\ge 1$) sao cho sau đúng $n$ bước chấm có thể quay về $O$. Chỉ ra một cách đi.
\item Sau đúng $8$ bước, chấm có thể quay về $O$ không?
\end{enumerate}
\end{probbox}
\probsol{
\textbf{Phân tích đề.} Như Bài 22, các bước luân phiên ngang -- dọc. Một nhóm phương gồm các bước có độ dài \emph{lẻ} $1,3,5,\dots$; nhóm kia gồm các bước độ dài \emph{chẵn} $2,4,6,\dots$ Chấm về $O$ khi và chỉ khi tổng có dấu của \emph{mỗi} nhóm bằng $0$.
\begin{recallbox}Quay về $O$ nghĩa là tổng dịch chuyển ngang bằng $0$ \emph{và} tổng dịch chuyển dọc bằng $0$. Đổi dấu một số hạng $a$ trong một tổng làm tổng thay đổi $2a$ -- một số chẵn -- nên \textbf{không làm đổi tính chẵn lẻ} của tổng.\end{recallbox}

\step{1}{Câu a}
Sau $10$ bước, nhóm lẻ gồm các độ dài $1,3,5,7,9$ (dù bước đầu là ngang hay dọc). Tổng có dấu $\pm1\pm3\pm5\pm7\pm9$ có cùng tính chẵn lẻ với $1+3+5+7+9=25$, tức là luôn \emph{lẻ}, nên không thể bằng $0$. Vậy theo phương đó chấm không về được vị trí ban đầu, tức chấm không ở $O$.

\step{2}{Câu b -- Loại các giá trị $n\le 6$}
Xét nhóm lẻ (theo một phương) với từng $n$:
\begin{center}
\begin{tabular}{|c|c|c|}
\hline
\rowcolor{tmBlue!10}$n$ & Nhóm độ dài lẻ & Vì sao tổng có dấu không thể bằng $0$\\ \hline
$1,\ 2$ & $\{1\}$ & $\pm 1\neq 0$\\ \hline
$3,\ 4$ & $\{1,3\}$ & $\pm1\pm3$ chỉ bằng $4, 2, -2, -4$\\ \hline
$5,\ 6$ & $\{1,3,5\}$ & tổng có dấu cùng tính chẵn lẻ với $1+3+5=9$ (lẻ)\\ \hline
\end{tabular}
\end{center}
Vậy với $n\le 6$, chấm không thể về $O$, nên $n\ge 7$.

\step{3}{Câu b -- Với $n=7$ thì được}
Nhóm lẻ: $\{1,3,5,7\}$, chọn $+1-3-5+7=0$. Nhóm chẵn: $\{2,4,6\}$, chọn $+2+4-6=0$.\par
Cho bước đầu là ngang. Cách đi: phải $1$, lên $2$, trái $3$, lên $4$, trái $5$, xuống $6$, phải $7$. Các vị trí lần lượt:
\[ (0;0)\to(1;0)\to(1;2)\to(-2;2)\to(-2;6)\to(-7;6)\to(-7;0)\to(0;0). \]
\begin{figC}\begin{scope}[scale=0.36]
\draw[Line,line width=0.3pt] (-8,-1) grid (2,7);
\draw[->,Navy,line width=0.5pt] (-8,0)--(2.4,0); \draw[->,Navy,line width=0.5pt] (0,-1)--(0,7.4);
\draw[tmBlue,line width=1.3pt,->] (0,0)--(1,0)--(1,2)--(-2,2)--(-2,6)--(-7,6)--(-7,0)--(-0.32,0);
\fill[AnsRed] (0,0) circle (3pt);
\node[font=\scriptsize,anchor=north east,inner sep=1.5pt] at (0,0) {$O$};
\foreach \p/\l/\a in {(0.5,0)/1/below,(1,1)/2/right,(-0.5,2)/3/below,(-2,4)/4/right,(-4.5,6)/5/above,(-7,3)/6/left,(-3.5,0)/7/below}{\node[\a,font=\scriptsize\bfseries,text=tmBlue!80!black] at \p {\l};}
\end{scope}\end{figC}
Vậy $n$ nhỏ nhất là $7$.

\step{4}{Câu c -- $n=8$}
Nhóm lẻ: $\{1,3,5,7\}$, vẫn chọn $+1-3-5+7=0$. Nhóm chẵn: $\{2,4,6,8\}$, chọn $+2-4-6+8=0$. Vậy \textbf{có thể}: phải $1$, lên $2$, trái $3$, xuống $4$, trái $5$, xuống $6$, phải $7$, lên $8$; các vị trí: $(1;0)$, $(1;2)$, $(-2;2)$, $(-2;-2)$, $(-7;-2)$, $(-7;-8)$, $(0;-8)$, $(0;0)$.
\begin{proofbox}Một tổng có dấu $\pm a_1\pm a_2\pm\dots$ bằng $0$ khi và chỉ khi ta chia được các số $a_1,a_2,\dots$ thành hai nhóm có tổng bằng nhau (nhóm mang dấu $+$ và nhóm mang dấu $-$). Ví dụ $\{1,3,5,7\}$: $1+7=3+5=8$. Vì vậy điều kiện cần là tổng tất cả các số phải \emph{chẵn}.\end{proofbox}
\begin{warnbox}Không được quên rằng \emph{cả hai} nhóm đều phải có tổng bằng $0$. Với $n=5$: nhóm chẵn $\{2,4\}$ không thể cho $0$, và nhóm lẻ $\{1,3,5\}$ cũng không -- chỉ cần một nhóm hỏng là đủ để loại.\end{warnbox}
\verify{Kiểm tra bằng máy tính mọi cách chọn chiều với $n=1,2,\dots,15$: chấm về được $O$ khi $n=7$, $n=8$ và $n=15$; mọi $n\le 6$ đều không được.}
\ans{a) Tổng có dấu của $1,3,5,7,9$ luôn lẻ nên khác $0$;\ b) $n=7$ (ví dụ: phải 1, lên 2, trái 3, lên 4, trái 5, xuống 6, phải 7);\ c) Có.}
}

% ------------------------------------------------------------------ 44
\begin{probbox}{44}{\stars{2}}
Một hình hộp chữ nhật kích thước $4\times 6\times n$ ($n$ là số nguyên dương) được ghép từ các hình lập phương đơn vị $1\times1\times1$. Gọi $A$ là diện tích toàn phần của hình hộp, $B$ là tổng diện tích toàn phần của tất cả các hình lập phương đơn vị.
\begin{enumerate}[label=\alph*),leftmargin=1.6em,itemsep=1pt]
\item Viết $A$, $B$ và tỉ số $\dfrac{B}{A}$ dưới dạng biểu thức của $n$.
\item Chứng minh rằng với mọi $n$, ta luôn có $\dfrac{B}{A}<7{,}2$.
\item Tìm tất cả các giá trị của $n$ để $\dfrac{B}{A}$ là số nguyên.
\end{enumerate}
\end{probbox}
\probsol{
\textbf{Phân tích đề.} Mở rộng của Bài 24, nhưng lần này mẹo ``tách phần nguyên'' không còn gọn như trước. Ta dùng một cách khác: chặn trên tỉ số, rồi thử từng giá trị nguyên có thể.
\begin{recallbox}Hình hộp $a\times b\times c$: diện tích toàn phần $2(ab+bc+ca)$, gồm $a\cdot b\cdot c$ hình lập phương đơn vị; mỗi hình lập phương đơn vị có diện tích toàn phần $6$.\end{recallbox}

\step{1}{Câu a}
$A=2(4\cdot 6+4\cdot n+6\cdot n)=2(24+10n)=48+20n$.\par
Số hình lập phương đơn vị: $4\cdot 6\cdot n=24n$; do đó $B=6\cdot 24n=144n$.\par
\[ \frac{B}{A}=\frac{144n}{48+20n}=\frac{4\cdot 36n}{4\cdot(12+5n)}=\frac{36n}{5n+12}. \]

\step{2}{Câu b}
Vì $5n+12>5n$, ta chia $36n$ cho một số \emph{lớn hơn} $5n$ nên được kết quả \emph{nhỏ hơn}:
\[ \frac{36n}{5n+12}<\frac{36n}{5n}=\frac{36}{5}=7{,}2. \]
\begin{recallbox}Với các số dương: cùng tử số, phân số nào có mẫu lớn hơn thì nhỏ hơn.\end{recallbox}

\step{3}{Câu c -- Đặt giá trị nguyên là $k$}
Nếu $\frac{B}{A}=k$ là số nguyên thì $k$ dương và $k<7{,}2$, nên $k\in\{1,2,3,4,5,6,7\}$. Ta có
\[ \frac{36n}{5n+12}=k \iff 36n=k(5n+12)=5kn+12k \iff 36n-5kn=12k \iff n(36-5k)=12k. \]
Vậy $n=\dfrac{12k}{36-5k}$, và ta chỉ cần thử $7$ giá trị của $k$.

\step{4}{Thử từng giá trị của $k$}
\begin{center}
\begin{tabular}{|c|c|c|c|c|c|c|c|}
\hline
\rowcolor{tmBlue!10}$k$ & $1$ & $2$ & $3$ & $4$ & $5$ & $6$ & $7$\\ \hline
$12k$ & $12$ & $24$ & $36$ & $48$ & $60$ & $72$ & $84$\\ \hline
$36-5k$ & $31$ & $26$ & $21$ & $16$ & $11$ & $6$ & $1$\\ \hline
$n=\frac{12k}{36-5k}$ & $\frac{12}{31}$ $\times$ & $\frac{24}{26}$ $\times$ & $\frac{36}{21}$ $\times$ & $\mathbf{3}$ & $\frac{60}{11}$ $\times$ & $\mathbf{12}$ & $\mathbf{84}$\\ \hline
\end{tabular}
\end{center}
(Dấu $\times$: không phải số nguyên.)
Vậy $n\in\{3;\,12;\,84\}$.
\begin{proofbox}Vì sao không bỏ sót? Mọi $n$ làm $\frac{B}{A}$ nguyên đều cho một $k$ nguyên trong khoảng $1$ đến $7$ (câu b), và với mỗi $k$ chỉ có đúng một $n$ thỏa $n(36-5k)=12k$. Ta đã thử hết $7$ khả năng.\end{proofbox}
\verify{
$n=3$: $A=48+60=108$, $B=432$, $\frac{B}{A}=4$ \checkmark.\quad
$n=12$: $A=48+240=288$, $B=1728$, $\frac{B}{A}=6$ \checkmark.\quad
$n=84$: $A=48+1680=1728$, $B=12\,096$, $\frac{B}{A}=7$ \checkmark.}
\begin{tipbox}Khi một phân số chứa $n$ không ``tách'' được gọn, hãy tìm xem nó luôn nhỏ hơn số nào. Nếu chỉ còn vài giá trị nguyên có thể, ta thử hết từng giá trị!\end{tipbox}
\ans{a) $A=48+20n$, $B=144n$, $\dfrac{B}{A}=\dfrac{36n}{5n+12}$;\ b) xem lời giải;\ c) $n=3$, $n=12$, $n=84$.}
}

% ------------------------------------------------------------------ 45
\begin{probbox}{45}{\stars{2}}
Hai tuyến xe buýt cắt nhau tại bến $X$. Xe A chạy đi chạy lại trên tuyến $P\text{--}X\text{--}S$, mỗi nhánh ($PX$ hoặc $XS$) mất đúng $4$ phút. Xe B chạy trên tuyến $Q\text{--}X\text{--}T$, mỗi nhánh mất đúng $6$ phút. Lúc $13$\,h\,$00$, xe A xuất phát từ $P$ và xe B xuất phát từ $Q$; các xe quay đầu ngay lập tức ở cuối tuyến và chạy liên tục.
\begin{enumerate}[label=\alph*),leftmargin=1.6em,itemsep=1pt]
\item Chứng minh rằng hai xe A và B không bao giờ cùng có mặt tại $X$.
\item Người ta cho thêm xe C chạy trên tuyến $R\text{--}X\text{--}U$, mỗi nhánh mất $10$ phút, xuất phát từ $R$ lúc $13$\,h\,$00$. Từ $13$\,h\,$00$ đến $19$\,h\,$00$, vào những thời điểm nào có hai xe cùng ở $X$?
\item Giả sử mỗi nhánh của tuyến xe C mất $m$ phút ($m$ là số nguyên dương bất kỳ). Chứng minh rằng xe C không thể vừa gặp xe A tại $X$ (vào một lúc nào đó), vừa gặp xe B tại $X$ (vào một lúc nào đó).
\end{enumerate}
\end{probbox}
\probsol{
\textbf{Phân tích đề.} Mở rộng của Bài 25. Đo thời gian $t$ bằng số phút kể từ $13$\,h\,$00$. Như ở Bài 25, một xe có nhánh dài $a$ phút sẽ tới $X$ lần đầu ở phút $a$, rồi cứ sau $2a$ phút (đi ra cuối tuyến và quay lại) lại tới $X$:
\[ t=a,\ 3a,\ 5a,\ 7a,\dots \quad\text{tức là } t=a\times(\text{số lẻ}). \]
Vậy xe A ở $X$ khi $t=4\times(\text{lẻ})$; xe B khi $t=6\times(\text{lẻ})$; xe C khi $t=m\times(\text{lẻ})$.
\begin{recallbox}Tích của hai số lẻ là số lẻ. Số chia hết cho $4$ thì khi chia cho $2$ được một số chẵn.\end{recallbox}

\step{1}{Câu a}
Thời điểm xe A ở $X$: $t=4\times(\text{lẻ})$ -- luôn \textbf{chia hết cho $4$}.\par
Thời điểm xe B ở $X$: $t=6\times(\text{lẻ})=2\times(3\times\text{lẻ})=2\times(\text{lẻ})$ -- chia cho $2$ được một số lẻ, nên \textbf{không chia hết cho $4$}.\par
Một số không thể vừa chia hết vừa không chia hết cho $4$, nên A và B không bao giờ cùng ở $X$.

\step{2}{Câu b -- xe C với $m=10$}
Xe C ở $X$ khi $t=10\times(\text{lẻ})=2\times(5\times\text{lẻ})$ -- không chia hết cho $4$, nên C \textbf{không bao giờ} gặp A tại $X$ (lý luận như câu a).\par
C gặp B khi $t=6\times(\text{lẻ})=10\times(\text{lẻ})$. Khi đó $\frac{t}{2}=3\times(\text{lẻ})=5\times(\text{lẻ})$ là số lẻ chia hết cho $3$ và $5$, tức là bội lẻ của $15$. Vậy $t=30\times(\text{lẻ})$. Ngược lại, $t=30\times(\text{lẻ})=6\times(5\times\text{lẻ})=10\times(3\times\text{lẻ})$ đúng là lúc cả hai xe ở $X$.\par
Từ $13$\,h\,$00$ đến $19$\,h\,$00$ là $0\le t\le 360$: $t=30,\ 90,\ 150,\ 210,\ 270,\ 330$.
\begin{center}
\begin{tabular}{|c|c|c|c|c|c|c|}
\hline
\rowcolor{tmBlue!10}$t$ (phút) & $30$ & $90$ & $150$ & $210$ & $270$ & $330$\\ \hline
Giờ & 13\,h\,30 & 14\,h\,30 & 15\,h\,30 & 16\,h\,30 & 17\,h\,30 & 18\,h\,30\\ \hline
\end{tabular}
\end{center}
Ở cả $6$ thời điểm này, hai xe cùng ở $X$ là B và C.

\step{3}{Câu c -- ``số lần chia đôi''}
Với mỗi số nguyên dương $t$, ta chia $t$ cho $2$ liên tiếp cho tới khi được số lẻ. Gọi số lần chia được là $h(t)$. Ví dụ: $12\to 6\to 3$ nên $h(12)=2$; $6\to3$ nên $h(6)=1$; $5$ đã lẻ nên $h(5)=0$.
\begin{proofbox}\textbf{Nhân với một số lẻ không làm thay đổi $h$.} Viết $m=2\cdot2\cdots2\times r$ (có $h(m)$ thừa số $2$, $r$ lẻ). Khi đó $m\times(\text{lẻ})=2\cdot2\cdots2\times(r\times\text{lẻ})$, trong đó $r\times\text{lẻ}$ vẫn là số lẻ. Vậy $m\times(\text{lẻ})$ cũng chia đôi được đúng $h(m)$ lần.\end{proofbox}
Áp dụng:
\begin{itemize}[leftmargin=1.3em,itemsep=1pt]
\item Mọi thời điểm xe A ở $X$: $h\big(4\times\text{lẻ}\big)=h(4)=2$.
\item Mọi thời điểm xe B ở $X$: $h\big(6\times\text{lẻ}\big)=h(6)=1$.
\item Mọi thời điểm xe C ở $X$: $h\big(m\times\text{lẻ}\big)=h(m)$.
\end{itemize}
Nếu C gặp A thì $h(m)=2$; nếu C gặp B thì $h(m)=1$. Một số $m$ không thể có cả $h(m)=2$ lẫn $h(m)=1$. Vậy C không thể gặp cả A lẫn B.
\begin{warnbox}Câu c không nói C \emph{luôn} không gặp ai: với $m=4$, C gặp A (lúc $t=4$); với $m=2$, C gặp B (lúc $t=6$). Điều ta chứng minh là C không thể gặp \emph{cả hai}.\end{warnbox}
\verify{Kiểm tra bằng máy tính với $m=1,2,\dots,12$: không có giá trị nào của $m$ để C gặp được cả A và B; $m=4$ và $m=12$ gặp A; $m=2,6,10$ gặp B.}
\ans{a) Thời điểm A ở $X$ chia hết cho $4$, thời điểm B ở $X$ thì không;\ b) $13$\,h\,$30$, $14$\,h\,$30$, $15$\,h\,$30$, $16$\,h\,$30$, $17$\,h\,$30$, $18$\,h\,$30$ (đều là xe B và C);\ c) xem lời giải.}
}

% ------------------------------------------------------------------ 46
\begin{probbox}{46}{\stars{1}}
Bạn Sam đếm tiến $7, 10, 13,\dots$ (mỗi lần thêm $3$) và dừng lại trước khi vượt quá $2026$. Bạn Tom đếm lùi $2026, 2021, 2016,\dots$ (mỗi lần bớt $5$) và dừng lại trước khi nhỏ hơn $7$.
\begin{enumerate}[label=\alph*),leftmargin=1.6em,itemsep=1pt]
\item Có bao nhiêu số có mặt trong cả hai dãy?
\item Tính tổng của tất cả các số ở câu a.
\item Bạn Zoe viết dãy các bội dương của $4$: $4, 8, 12,\dots$ Có bao nhiêu số có mặt trong cả ba dãy của Sam, Tom và Zoe? Số lớn nhất trong đó là số nào?
\end{enumerate}
\end{probbox}
\probsol{
\textbf{Phân tích đề.} Mở rộng của Bài 26. Như ở Bài 26: một số thuộc dãy của Sam $\iff$ nó chia $3$ dư $1$ và $7\le N\le 2026$ (vì $7=3\cdot 2+1$, mỗi lần cộng $3$ thì số dư khi chia cho $3$ không đổi); một số thuộc dãy của Tom $\iff$ nó chia $5$ dư $1$ và $7\le N\le 2026$ (vì $2026=5\cdot 405+1$). Bây giờ ta tìm \emph{tất cả} các số chung.

\step{1}{Câu a -- Dấu hiệu của số chung}
Số $N$ chung $\iff$ $N$ chia $3$ dư $1$ và chia $5$ dư $1$ $\iff$ $N-1$ chia hết cho $3$ và cho $5$ $\iff$ $N-1$ chia hết cho $15$. Vậy $N=15j+1$.
\begin{recallbox}Một số chia hết cho $3$ và cho $5$ (hai số nguyên tố cùng nhau) thì chia hết cho $3\cdot 5=15$.\end{recallbox}

\step{2}{Câu a -- Đếm}
Cần $7\le 15j+1\le 2026$. Số nhỏ nhất: $j=1$ cho $16$ ($j=0$ cho $1<7$, loại). Số lớn nhất: $15j+1\le 2026 \Rightarrow 15j\le 2025 \Rightarrow j\le 135$; $j=135$ cho $2026$.\par
Vậy các số chung là $16, 31, 46,\dots, 2026$ ứng với $j=1,2,\dots,135$: có \textbf{$135$ số}.

\step{3}{Câu b -- Tổng (cách ghép cặp của Gauss)}
\begin{recallbox}Tổng của một dãy số cách đều $=$ (số đầu $+$ số cuối) $\times$ (số số hạng) $:2$. Lý do: viết dãy theo thứ tự xuôi và ngược rồi cộng từng cặp, mỗi cặp đều bằng số đầu $+$ số cuối.\end{recallbox}
\[ 16+31+\dots+2026=\frac{(16+2026)\cdot 135}{2}=\frac{2042\cdot 135}{2}=1021\cdot 135=137\,835. \]

\step{4}{Câu c -- Thêm điều kiện chia hết cho $4$}
Xét các số chung $16, 31, 46, 61, 76, 91, 106, 121, 136,\dots$ Số nào chia hết cho $4$?
\[ 16\ \checkmark,\ 31,\ 46,\ 61,\ 76\ \checkmark,\ 91,\ 106,\ 121,\ 136\ \checkmark,\dots \]
Cứ $4$ số liên tiếp thì có đúng $1$ số chia hết cho $4$, và các số đó cách nhau $4\cdot 15=60$: $16, 76, 136, \dots$
\begin{proofbox}Mỗi lần đi tiếp trong dãy, số tăng $15$; tăng $15$ thì số dư khi chia $4$ tăng $3$ (vì $15=12+3$). Các số dư khi chia $4$ của $16,31,46,61$ lần lượt là $0,3,2,1$, rồi lặp lại. Vậy chính xác cứ $4$ số có một số chia hết cho $4$, và hai số liền nhau như vậy cách nhau $4\times15=60$.\end{proofbox}
Các số chung của cả ba dãy: $16+60i$ với $16+60i\le 2026$, tức $60i\le 2010$, $i\le 33{,}5$, nên $i=0,1,\dots,33$: có \textbf{$34$ số}. Số lớn nhất: $16+60\cdot 33=16+1980=1996$.
\verify{$1996=4\cdot 499$ \checkmark; $1996-7=1989=3\cdot 663$ \checkmark; $2026-1996=30=5\cdot 6$ \checkmark. Số tiếp theo $1996+60=2056>2026$, bị loại.}
\ans{a) $135$ số;\ b) $137\,835$;\ c) $34$ số, lớn nhất là $1996$.}
}

% ------------------------------------------------------------------ 47
\begin{probbox}{47}{\stars{2}}
Emma và Noah cùng xuất phát tại điểm $A$ trên một đường chạy hình tròn, mỗi người chạy với vận tốc không đổi.
\begin{enumerate}[label=\alph*),leftmargin=1.6em,itemsep=1pt]
\item Hai bạn chạy ngược chiều nhau, tỉ số vận tốc Emma : Noah là $5:3$. Hai bạn gặp nhau ở bao nhiêu điểm khác nhau trên đường chạy?
\item Nếu hai bạn chạy \emph{cùng chiều} (vẫn với tỉ số $5:3$) thì Emma vượt Noah ở bao nhiêu điểm khác nhau?
\item Tìm tất cả các tỉ số vận tốc $p:q$ (với $p>q$, $p$ và $q$ là các số nguyên dương nguyên tố cùng nhau) sao cho khi chạy ngược chiều nhau, hai bạn gặp nhau ở đúng $10$ điểm khác nhau.
\end{enumerate}
\end{probbox}
\probsol{
\textbf{Phân tích đề.} Mở rộng của Bài 27. Ta dùng lại hai ý: (1) chạy ngược chiều thì giữa hai lần gặp liên tiếp, tổng quãng đường của hai bạn bằng $1$ vòng; (2) quãng đường tỉ lệ với vận tốc.

\step{1}{Câu a -- Chia vòng thành 8 đoạn}
Tỉ số $5:3$ nên ta chia vòng chạy thành $5+3=8$ đoạn bằng nhau, đánh số các điểm chia $0,1,2,\dots,7$ theo chiều chạy của Emma (điểm $0$ là $A$). Giữa hai lần gặp, hai bạn cùng chạy $8$ đoạn, trong đó Emma chạy $5$ đoạn, Noah $3$ đoạn.\par
Vị trí lần gặp thứ $k$ là điểm cách $A$ đúng $5k$ đoạn; cứ đủ $8$ đoạn là một vòng nên ta bớt đi các bội của $8$:
\begin{center}\small
\begin{tabular}{|c|c|c|c|c|c|c|c|c|c|}
\hline
\rowcolor{tmBlue!10}Lần gặp $k$ & $1$ & $2$ & $3$ & $4$ & $5$ & $6$ & $7$ & $8$ & $9$\\ \hline
$5k$ & $5$ & $10$ & $15$ & $20$ & $25$ & $30$ & $35$ & $40$ & $45$\\ \hline
Điểm (bớt bội của $8$) & $5$ & $2$ & $7$ & $4$ & $1$ & $6$ & $3$ & $0$ & $5$\\ \hline
\end{tabular}
\end{center}
Tám lần gặp đầu rơi vào đủ $8$ điểm khác nhau $0,1,\dots,7$; lần thứ $9$ trở lại điểm $5$ và mọi thứ lặp lại. Vậy có \textbf{$8$ điểm gặp}.
\begin{figC}
\draw[Ink,line width=1.1pt] (0,0) circle (1.3);
\foreach \i/\k in {0/8,1/5,2/2,3/7,4/4,5/1,6/6,7/3}{
  \pgfmathsetmacro\rvabAng{90+45*\i}
  \fill[AnsRed] (\rvabAng:1.3) circle (2.6pt);
  \node[font=\scriptsize,anchor=\rvabAng+180,inner sep=2pt] at (\rvabAng:1.42) {điểm \i};
  \node[font=\scriptsize\bfseries,text=tmBlue] at (\rvabAng:0.98) {\k};}
\node[font=\scriptsize,text=tmBlue,align=center] at (0,0) {lần\\gặp};
\end{figC}
(Số xanh bên trong: lần gặp thứ mấy rơi vào điểm đó; các điểm được đánh số theo chiều chạy của Emma.)

\step{2}{Câu b -- Chạy cùng chiều}
Chạy cùng chiều thì Emma vượt Noah mỗi khi Emma chạy \emph{hơn} Noah đúng $1$ vòng. Trong cùng thời gian, Emma chạy $5$ phần thì Noah chạy $3$ phần, Emma hơn $2$ phần. Vậy ``hơn $1$ vòng'' ứng với $2$ phần $=1$ vòng, tức mỗi phần là $\frac12$ vòng, và Emma đã chạy $5\cdot\frac12=2\frac12$ vòng.\par
Lần vượt thứ nhất: Emma chạy $2\frac12$ vòng, ở điểm đối diện $A$. Lần thứ hai: $5$ vòng, đúng tại $A$. Rồi lặp lại. Vậy có \textbf{$2$ điểm} (điểm $A$ và điểm đối diện với $A$).

\step{3}{Câu c -- Quy luật tổng quát}
Làm như câu a với tỉ số $p:q$: chia vòng thành $p+q$ đoạn; lần gặp thứ $k$ ở điểm cách $A$ đúng $k\cdot p$ đoạn (bớt các bội của $p+q$).
\begin{proofbox}\textbf{Nếu $p$, $q$ nguyên tố cùng nhau thì có đúng $p+q$ điểm gặp.}\par
\emph{Lần thứ $p+q$ trở về $A$:} vì $(p+q)\cdot p$ là bội của $p+q$.\par
\emph{Các lần $1,2,\dots,p+q$ ở các điểm khác nhau:} giả sử lần $i$ và lần $j$ ($i<j\le p+q$) trùng điểm. Khi đó $j\cdot p-i\cdot p=(j-i)\cdot p$ là bội của $p+q$. Mọi ước chung của $p$ và $p+q$ cũng là ước của $(p+q)-p=q$, mà $p,q$ nguyên tố cùng nhau, nên $p$ và $p+q$ không có ước nguyên tố chung. Do đó mọi thừa số nguyên tố của $p+q$ phải nằm trong $j-i$, tức $j-i$ chia hết cho $p+q$. Nhưng $0<j-i<p+q$ -- vô lý.\par
Vậy có đúng $p+q$ điểm gặp, cách đều nhau trên vòng tròn.\end{proofbox}
Cần $p+q=10$ với $p>q\ge 1$ nguyên tố cùng nhau:
\begin{center}
\begin{tabular}{|c|c|c|c|c|}
\hline
\rowcolor{tmBlue!10}$p:q$ & $9:1$ & $8:2$ & $7:3$ & $6:4$\\ \hline
Nguyên tố cùng nhau? & có & không ($=4:1$) & có & không ($=3:2$)\\ \hline
Số điểm gặp & $10$ & $5$ & $10$ & $5$\\ \hline
\end{tabular}
\end{center}
($5:5$ bị loại vì $p>q$.) Vậy có hai tỉ số: \textbf{$9:1$ và $7:3$}.
\begin{warnbox}Tỉ số $8:2$ có $8+2=10$ nhưng thực chất là $4:1$, chỉ cho $4+1=5$ điểm gặp. Luôn rút gọn tỉ số trước khi áp dụng quy luật ``$p+q$ điểm''.\end{warnbox}
\verify{Kiểm tra lại Bài 27 bằng quy luật: tỉ số $2:1$ cho $2+1=3$ điểm \checkmark. Mô phỏng bằng máy tính với các tỉ số $p:q$ tối giản, $p\le 19$: chỉ có $9:1$ và $7:3$ cho $10$ điểm.}
\ans{a) $8$ điểm;\ b) $2$ điểm;\ c) $9:1$ và $7:3$.}
}

% ------------------------------------------------------------------ 48
\begin{probbox}{48}{\stars{2}}
Ta gọi một dãy số nguyên dương \emph{tăng} là \emph{dãy Fibo} nếu kể từ số thứ ba, mỗi số bằng tổng của hai số đứng ngay trước nó (ví dụ: $11, 20, 31, 51, 82$ như ở Bài 28).
\begin{enumerate}[label=\alph*),leftmargin=1.6em,itemsep=1pt]
\item Tìm tất cả các dãy Fibo gồm $6$ số có số cuối cùng bằng $100$.
\item Có bao nhiêu dãy Fibo gồm $5$ số có số cuối cùng bằng $200$?
\item Chứng minh rằng không có dãy Fibo nào gồm $6$ số có số cuối cùng bằng $30$.
\end{enumerate}
\end{probbox}
\probsol{
\textbf{Phân tích đề.} Như Bài 28, cả dãy được xác định bởi hai số đầu $a$ và $b$; dãy tăng khi và chỉ khi $a<b$.

\step{1}{Biểu diễn các số theo $a$ và $b$}
\begin{center}\small
\begin{tabular}{|c|c|c|c|c|c|}
\hline
\rowcolor{tmBlue!10}Số thứ 1 & Số thứ 2 & Số thứ 3 & Số thứ 4 & Số thứ 5 & Số thứ 6\\ \hline
$a$ & $b$ & $a+b$ & $a+2b$ & $2a+3b$ & $(a+2b)+(2a+3b)=3a+5b$\\ \hline
\end{tabular}
\end{center}
\begin{tipbox}Các hệ số $1,1,2,3,5,8,\dots$ xuất hiện trong bảng chính là dãy Fibonacci nổi tiếng!\end{tipbox}

\step{2}{Câu a -- Phương trình $3a+5b=100$}
Điều kiện: $a\ge 1$, $a<b$, $3a+5b=100$.
\begin{itemize}[leftmargin=1.3em,itemsep=1pt]
\item Từ $a\ge 1$: $5b=100-3a\le 97$, nên $b\le 19$.
\item Từ $a<b$: $100=3a+5b<3b+5b=8b$, nên $b>12{,}5$, tức $b\ge 13$.
\item $3a=100-5b$ phải chia hết cho $3$.
\end{itemize}
Thử $b$ từ $13$ đến $19$:
\begin{center}\small
\begin{tabular}{|c|c|c|c|c|c|c|c|}
\hline
\rowcolor{tmBlue!10}$b$ & $13$ & $14$ & $15$ & $16$ & $17$ & $18$ & $19$\\ \hline
$100-5b$ & $35$ & $30$ & $25$ & $20$ & $15$ & $10$ & $5$\\ \hline
$a=(100-5b):3$ & $\times$ & $\mathbf{10}$ & $\times$ & $\times$ & $\mathbf{5}$ & $\times$ & $\times$\\ \hline
\end{tabular}
\end{center}
Có hai dãy: $\mathbf{10, 14, 24, 38, 62, 100}$ và $\mathbf{5, 17, 22, 39, 61, 100}$.

\step{3}{Câu b -- Phương trình $2a+3b=200$}
Vì $2a$ và $200$ chẵn nên $3b$ chẵn, suy ra $b$ chẵn: $b=2t$. Khi đó $2a=200-6t$, tức $a=100-3t$.
\begin{itemize}[leftmargin=1.3em,itemsep=1pt]
\item $a\ge 1$: $100-3t\ge 1 \Rightarrow t\le 33$.
\item $a<b$: $100-3t<2t \Rightarrow 100<5t \Rightarrow t>20 \Rightarrow t\ge 21$.
\end{itemize}
Vậy $t=21,22,\dots,33$: có $33-21+1=13$ dãy. (Ví dụ $t=21$: $37, 42, 79, 121, 200$; $t=33$: $1, 66, 67, 133, 200$.)
\begin{warnbox}Với $t=20$ ta được $a=b=40$: dãy $40,40,80,120,200$ \emph{không tăng} ở bước đầu, nên bị loại. Đếm số nguyên từ $21$ đến $33$ là $33-21+1=13$, không phải $12$.\end{warnbox}

\step{4}{Câu c -- Không có dãy nào kết thúc bằng $30$}
Cần $3a+5b=30$ với $1\le a<b$.
\begin{itemize}[leftmargin=1.3em,itemsep=1pt]
\item $3a=30-5b\ge 3$ nên $5b\le 27$, tức $b\le 5$.
\item $5b=30-3a$ chia hết cho $3$, nên $b$ chia hết cho $3$; với $b\le 5$ chỉ có $b=3$.
\item $b=3$ cho $3a=15$, $a=5$. Nhưng $a=5>3=b$ -- dãy $5,3,8,11,19,30$ không tăng.
\end{itemize}
Vậy không có dãy Fibo $6$ số nào kết thúc bằng $30$.
\begin{proofbox}Có thể lập luận gọn hơn: vì $a<b$ nên $a\ge 1$, $b\ge 2$, và số thứ sáu $3a+5b<3b+5b=8b$. Muốn có $3a+5b=30$ thì phải có $b>30:8$, tức $b\ge 4$; cùng với $b\le 5$ và $b$ chia hết cho $3$ thì không còn giá trị nào.\end{proofbox}
\verify{Thử bằng máy tính mọi cặp $1\le a<b$: đúng $2$ dãy ở câu a, $13$ dãy ở câu b và $0$ dãy ở câu c.}
\ans{a) $10,14,24,38,62,100$ và $5,17,22,39,61,100$;\ b) $13$ dãy;\ c) xem lời giải.}
}

% ------------------------------------------------------------------ 49
\begin{probbox}{49}{\stars{1}}
Bốn bạn Anna, Ben, Chloe và David, mỗi bạn có ít nhất một viên bi. Biết rằng:
\begin{itemize}[leftmargin=1.6em,itemsep=0pt]
\item[(1)] Số bi của Anna và Ben cộng lại bằng số bi của Chloe và David cộng lại.
\item[(2)] Số bi của Anna và David cộng lại nhiều hơn số bi của Ben và Chloe cộng lại.
\item[(3)] Ben có nhiều bi hơn Anna và Chloe cộng lại.
\end{itemize}
\begin{enumerate}[label=\alph*),leftmargin=1.6em,itemsep=1pt]
\item Sắp xếp bốn bạn theo thứ tự số bi từ ít đến nhiều.
\item Biết cả bốn bạn có tổng cộng $30$ viên bi. Chloe có nhiều nhất bao nhiêu viên bi? Khi đó mỗi bạn có bao nhiêu viên?
\end{enumerate}
\end{probbox}
\probsol{
\textbf{Phân tích đề.} Mở rộng của Bài 29: thay vì so sánh trực tiếp từng cặp, ở đây các thông tin là so sánh các \emph{tổng}. Gọi số bi của Anna, Ben, Chloe, David lần lượt là $a, b, c, d$ (số nguyên dương). Ta có:
\[ \text{(1) } a+b=c+d,\qquad \text{(2) } a+d>b+c,\qquad \text{(3) } b>a+c. \]
\begin{recallbox}Cộng (hoặc trừ) hai vế của một bất đẳng thức với hai vế tương ứng của một đẳng thức, ta được một bất đẳng thức cùng chiều: nếu $X>Y$ và $M=N$ thì $X+M>Y+N$ và $X-M>Y-N$.\end{recallbox}

\step{1}{So sánh Anna và Chloe: cộng (2) với (1)}
$(a+d)+(a+b)>(b+c)+(c+d)$, tức $2a+b+d>b+d+2c$. Bớt $b+d$ ở hai vế: $2a>2c$, vậy $a>c$: \textbf{Anna nhiều hơn Chloe}.

\step{2}{So sánh Ben và David: lấy (2) trừ (1)}
$(a+d)-(a+b)>(b+c)-(c+d)$, tức $d-b>b-d$. Cộng $b+d$ vào hai vế: $2d>2b$, vậy $d>b$: \textbf{David nhiều hơn Ben}.

\step{3}{So sánh Anna và Ben: dùng (3)}
$b>a+c\ge a+1>a$ (vì $c\ge 1$). Vậy \textbf{Ben nhiều hơn Anna}.

\step{4}{Câu a -- Kết luận}
Ghép ba kết quả lại: $c<a$, \ $a<b$, \ $b<d$, nên $c<a<b<d$.\par
Thứ tự từ ít đến nhiều: \textbf{Chloe, Anna, Ben, David}.

\step{5}{Câu b -- Dùng tổng bằng 30}
Từ (1), $a+b=c+d$ mà $(a+b)+(c+d)=30$, nên $a+b=c+d=15$. Suy ra $b=15-a$.\par
Thay vào (3): $15-a>a+c$, tức $2a+c<15$; vì $2a+c$ là số nguyên nên $2a+c\le 14$.\par
Theo câu a, $a>c$ nên $a\ge c+1$. Do đó $2(c+1)+c\le 2a+c\le 14$, tức $3c+2\le 14$, $3c\le 12$, $c\le 4$.

\step{6}{Câu b -- Kiểm tra $c=4$ có xảy ra không}
Với $c=4$: cần $a\ge 5$ và $2a+4\le 14$, tức $a\le 5$. Vậy $a=5$, $b=15-5=10$, $d=15-4=11$.\par
Kiểm tra cả ba điều kiện:
(1) $5+10=15=4+11$ \checkmark;\quad (2) $5+11=16>14=10+4$ \checkmark;\quad (3) $10>5+4=9$ \checkmark.
\begin{warnbox}Chứng minh ``$c\le 4$'' mới chỉ là một nửa: phải chỉ ra một trường hợp cụ thể có $c=4$ thì mới khẳng định được $4$ là giá trị lớn nhất.\end{warnbox}
\verify{Liệt kê bằng máy tính mọi bộ $(a;b;c;d)$ có tổng $30$ thỏa mãn (1), (2), (3): có $12$ bộ, trong đó $c$ lớn nhất bằng $4$, chỉ ở bộ $(5;10;4;11)$; và trong mọi bộ đều có $c<a<b<d$.}
\ans{a) Chloe, Anna, Ben, David;\ b) Chloe có nhiều nhất $4$ viên; khi đó Anna có $5$, Ben $10$, Chloe $4$, David $11$ viên.}
}

% ------------------------------------------------------------------ 50
\begin{probbox}{50}{\stars{1}}
Một cửa hàng bán kẹo loại A giá \$12 một ki-lô-gam và kẹo loại B giá \$20 một ki-lô-gam.
\begin{enumerate}[label=\alph*),leftmargin=1.6em,itemsep=1pt]
\item Cần trộn bao nhiêu gam mỗi loại để được $1$~kg kẹo hỗn hợp có giá \$15 một ki-lô-gam (giá trị của hỗn hợp bằng đúng tổng giá trị kẹo đem trộn)?
\item Phải thêm bao nhiêu gam kẹo loại C giá \$30 một ki-lô-gam vào $1$~kg hỗn hợp ở câu a để được hỗn hợp mới có giá \$18 một ki-lô-gam?
\item Chứng minh rằng: trộn kẹo A với kẹo B (có mặt cả hai loại) theo bất kỳ tỉ lệ nào, giá của hỗn hợp luôn lớn hơn \$12 và nhỏ hơn \$20 một ki-lô-gam.
\end{enumerate}
\end{probbox}
\probsol{
\textbf{Phân tích đề.} Mở rộng của Bài 30. Nguyên tắc: \emph{tổng số tiền của hỗn hợp $=$ tổng số tiền của các phần đem trộn}. Ta đo khối lượng bằng ki-lô-gam và số tiền bằng đô la.
\begin{recallbox}Số tiền $=$ giá (cho $1$~kg) $\times$ khối lượng (kg). Giá trung bình của hỗn hợp $=$ tổng số tiền $:$ tổng khối lượng.\end{recallbox}

\step{1}{Câu a -- Gọi ẩn và lập phương trình}
Gọi khối lượng kẹo A là $x$ (kg), với $0<x<1$; khi đó khối lượng kẹo B là $1-x$ (kg).\par
Tiền kẹo A: $12x$; tiền kẹo B: $20(1-x)=20-20x$; tiền của $1$~kg hỗn hợp: $15\cdot 1=15$ (đô la). Ta có phương trình
\[ 12x+20-20x=15 \iff 20-8x=15 \iff 8x=5 \iff x=\frac{5}{8}=0{,}625. \]
Vậy cần $0{,}625$~kg $=625$~g kẹo A và $1000-625=375$~g kẹo B.
\verify{$12\cdot 0{,}625+20\cdot 0{,}375=7{,}5+7{,}5=15$ (đô la) cho $1$~kg, đúng giá \$15 một ki-lô-gam \checkmark.}

\step{2}{Câu a -- Cách khác: so sánh với giá \$15 trên trục số}
\begin{figC}
\draw[->,Navy,line width=1pt] (-0.5,0)--(5.4,0);
\foreach \x/\l in {0/12,1.8/15,4.8/20}{\draw[Navy,line width=1pt] (\x,0.12)--(\x,-0.12); \node[below,font=\small] at (\x,-0.12) {\l};}
\node[above,font=\small\bfseries,text=tmBlue] at (0,0.15) {A};
\node[above,font=\small\bfseries,text=AnsRed] at (1.8,0.15) {hỗn hợp};
\node[above,font=\small\bfseries,text=tmOrange!85!black] at (4.8,0.15) {B};
\node[anchor=west,font=\footnotesize,text=Slate] at (5.45,0) {đô la/kg};
\draw[Slate,dashed,line width=0.5pt] (0,-0.6)--(0,-1.05) (1.8,-0.6)--(1.8,-1.05) (4.8,-0.6)--(4.8,-1.05);
\draw[<->,tmBlue!72!black,line width=0.7pt] (0,-0.9)--(1.8,-0.9);
\draw[<->,tmBlue!72!black,line width=0.7pt] (1.8,-0.9)--(4.8,-0.9);
\node[below,font=\footnotesize] at (0.9,-0.95) {$15-12=3$};
\node[below,font=\footnotesize] at (3.3,-0.95) {$20-15=5$};
\end{figC}
Mỗi ki-lô-gam kẹo A rẻ hơn giá \$15 là \$3; mỗi ki-lô-gam kẹo B đắt hơn giá \$15 là \$5. Phần ``rẻ hơn'' phải bù đúng phần ``đắt hơn'': $3x=5(1-x)$. Giải: $3x=5-5x$, $8x=5$, $x=0{,}625$ -- trùng với Bước 1. Từ $3x=5y$ (với $y=1-x$) suy ra $\dfrac{x}{y}=\dfrac{5}{3}$: khối lượng A và B tỉ lệ với $5$ và $3$, tức là tỉ lệ \emph{nghịch đảo} với hai khoảng cách $3$ và $5$.

\step{3}{Câu b}
Gọi khối lượng kẹo C thêm vào là $y$ (kg), $y>0$. Tổng số tiền: $15+30y$; tổng khối lượng: $1+y$ (kg). Muốn giá \$18 một ki-lô-gam thì
\[ 15+30y=18(1+y)=18+18y \iff 30y-18y=18-15 \iff 12y=3 \iff y=0{,}25. \]
Vậy cần thêm $0{,}25$~kg $=250$~g kẹo C.
\verify{Tổng số tiền $15+30\cdot 0{,}25=15+7{,}5=22{,}5$ (đô la) cho $1{,}25$~kg: $22{,}5:1{,}25=18$ \checkmark. (Theo cách ở Bước 2: $18$ cách $15$ là $3$, cách $30$ là $12$, nên $3\cdot 1=12y$, $y=0{,}25$ \checkmark.)}

\step{4}{Câu c}
Gọi khối lượng kẹo A là $x$ (kg), kẹo B là $y$ (kg), với $x>0$, $y>0$. Giá hỗn hợp là $G=\dfrac{12x+20y}{x+y}$ (đô la/kg).
\begin{itemize}[leftmargin=1.3em,itemsep=1pt]
\item Vì $y>0$ nên $20y>12y$, do đó $12x+20y>12x+12y=12(x+y)$. Chia hai vế cho số dương $x+y$: $G>12$.
\item Vì $x>0$ nên $12x<20x$, do đó $12x+20y<20x+20y=20(x+y)$. Chia hai vế cho $x+y$: $G<20$.
\end{itemize}
Vậy $12<G<20$: giá hỗn hợp luôn nằm \emph{giữa} giá của hai loại kẹo.
\begin{warnbox}Không thể trộn A và B để được giá \$25 hay \$10 một ki-lô-gam: giá trung bình không bao giờ vượt ra ngoài khoảng từ giá thấp nhất đến giá cao nhất.\end{warnbox}
\begin{parentbox}Có thể minh họa câu c bằng việc pha nước: trộn nước $20^\circ$C với nước $60^\circ$C thì được nước ấm ở giữa, không bao giờ nóng hơn $60^\circ$C. ``Trung bình luôn nằm ở giữa'' là một ý tưởng rất hữu ích.\end{parentbox}
\ans{a) $625$~g kẹo A và $375$~g kẹo B;\ b) $250$~g kẹo C;\ c) xem lời giải.}
}
